%%%%%%%%%%%%%%%%%%%%%%%%%%%%%%%%%%%%%%%%%
% FRI NLP Report
%%%%%%%%%%%%%%%%%%%%%%%%%%%%%%%%%%%%%%%%%

\documentclass[fleqn,moreauthors,10pt]{ds_report}
\usepackage[english]{babel}
\usepackage{hyperref}
\graphicspath{{fig/}}

\JournalInfo{FRI Natural Language Processing Course 2026}
\Archive{Project report}

\PaperTitle{
\begin{center}
Report 2\\[0.3em]
Domain-Specific Fitness NLP Assistant Using Retrieval-Augmented Generation
\end{center}
}
\Authors{Domen Pahole (Team leader), Luka Rizman, Nina Trivić}

\affiliation{\textit{Advisors: Aleš Žagar}}

\Keywords{Natural Language Processing, Fitness NLP, Retrieval-Augmented Generation, Semantic Retrieval, Domain Adaptation}
\newcommand{\keywordname}{Keywords}

\Abstract{
This project develops a domain-specific NLP assistant for fitness applications. 
The system leverages structured exercise data to enable semantic retrieval and contextual response generation. 
This report presents the first implementation results and a comparison between a baseline LLM and a RAG-based approach.
}

\begin{document}
\flushbottom
\maketitle
\thispagestyle{empty}

\section*{Introduction}

In this report, we will introduce first results of testing our own RAG model in comparison with Llama 3.1 8B model. The topic is Fitness Assistant. 

\section*{Datasets for Current Testing}

For the current phase of testing, we use a single dataset, \texttt{MegaGymDataset.csv}. 
The dataset contains 2,918 exercise entries with structured metadata describing each exercise.

Each entry includes the following attributes:

\begin{itemize}
\item \textbf{Title} – name of the exercise
\item \textbf{Desc} – textual description of how the exercise is performed
\item \textbf{Type} – category of exercise (e.g. Strength, Cardio)
\item \textbf{BodyPart} – primary muscle group targeted
\item \textbf{Equipment} – required equipment
\item \textbf{Level} – difficulty level (e.g. Beginner, Intermediate, Expert)
\item \textbf{Rating} – numerical user rating
\item \textbf{RatingDesc} – textual interpretation of the rating
\end{itemize}

An example entry from the dataset:

\begin{itemize}
\item \textbf{Title}: Dumbbell chest fly
\item \textbf{Desc}: Description of the exercise execution
\item \textbf{Type}: Strength
\item \textbf{BodyPart}: Chest
\item \textbf{Equipment}: Dumbbell
\item \textbf{Level}: Intermediate
\end{itemize}

\section*{Results}

In this section we present the first evaluation of our prototype RAG system. 
The goal is not a full-scale scientific evaluation, but rather to verify whether the core components already provide meaningful improvements over a baseline LLM setup.

We evaluate two aspects separately:

\begin{itemize}
\item Retrieval quality (vector search in Qdrant)
\item End-to-end answer quality (Baseline LLM vs RAG)
\end{itemize}

\section*{Retrieval Evaluation}

We first evaluate only the retrieval component, without using the LLM. The goal is to check whether the system can find relevant exercises for a given query.

Each test query is annotated with expected properties such as:

\begin{itemize}
\item target body part
\item equipment
\item optional difficulty level
\item expected keywords
\end{itemize}

The retrieval pipeline is:

\begin{enumerate}
\item question $\rightarrow$ embedding (MiniLM)
\item embedding $\rightarrow$ Qdrant
\item top-5 most similar exercises returned
\item evaluation against expected criteria
\end{enumerate}

\subsection*{Metrics}

We use standard information retrieval metrics:

\begin{itemize}
\item Hit@k (whether a relevant result appears in top k)
\item Precision@5
\item Mean Reciprocal Rank (MRR)
\end{itemize}

\subsection*{Results}

\begin{center}
\begin{tabular}{l r}
\hline
Metric & Value \\
\hline
Number of queries & 12 \\
Hit@1 & 0.500 \\
Hit@3 & 0.917 \\
Hit@5 & 0.917 \\
Precision@5 & 0.583 \\
MRR & 0.694 \\
\hline
\end{tabular}
\end{center}

\subsection*{Interpretation}

The results show that the retrieval component already performs reasonably well.

Hit@5 = 0.917 means that in 11 out of 12 queries, at least one relevant exercise appears in the top 5 results. This is sufficient for a RAG system, where the LLM can use multiple candidates as context.

However, Hit@1 = 0.500 indicates that the top result is not always optimal. This is expected at this stage, as we use a general-purpose embedding model and datasets with inconsistent metadata.

A notable failure case is a query for beginner bodyweight leg exercises, where retrieval fails to return relevant results. This highlights limitations of:

\begin{itemize}
\item missing or inconsistent dataset fields
\item lack of filtering by metadata (e.g. equipment)
\end{itemize}

Overall, retrieval is already usable, but clearly the main bottleneck for further improvement.

\section*{Baseline LLM vs RAG}

We compare two approaches using the same model (Llama 3.1 8B via Ollama):

\begin{itemize}
\item Baseline: direct question $\rightarrow$ LLM
\item RAG: question $\rightarrow$ retrieval $\rightarrow$ LLM with context
\end{itemize}

The difference lies only in the prompt, not in the model.

\subsection*{Key Observations}

\begin{center}
\begin{tabular}{l l l}
\hline
Aspect & Baseline & RAG \\
\hline
Knowledge source & General & Dataset + LLM \\
Exercise names & Generic & More specific \\
Grounding & None & Dataset-grounded \\
Hallucination risk & Higher & Lower (if retrieval good) \\
Weakness & Not dataset-aligned & Depends on retrieval \\
\hline
\end{tabular}
\end{center}

\subsection*{Example Insights}

For queries such as:

\textit{"What chest exercises can I do with dumbbells?"}

the baseline produces reasonable general answers, but may include exercises not present in our dataset.

The RAG system, on the other hand, generates answers based on retrieved entries, leading to:

\begin{itemize}
\item more consistent exercise naming
\item better alignment with the dataset
\item structured explanations based on retrieved metadata
\end{itemize}

However, in cases where retrieval fails (e.g. broad queries like beginner bodyweight legs), the RAG answer becomes worse than the baseline. This confirms that RAG performance is highly dependent on retrieval quality.

\subsection*{Discussion}

The comparison highlights an important insight:

\begin{itemize}
\item Baseline LLM is strong in general reasoning and common knowledge
\item RAG is strong in grounding answers to a specific dataset
\item The overall system quality is dominated by retrieval performance
\end{itemize}

This suggests that the next development focus should not be a larger model, but improvements in retrieval:

\begin{itemize}
\item metadata filtering (BodyPart, Equipment, Level)
\item dataset cleaning and deduplication
\item improved evaluation dataset
\end{itemize}

\section*{Conclusion of Current Phase}

At this stage, we have successfully implemented a working RAG prototype:

\begin{itemize}
\item vector-based retrieval using Qdrant
\item embedding model integration
\item prompt-based grounding for LLM
\item initial evaluation framework
\end{itemize}

The system already demonstrates the core advantage of RAG: grounding responses in a controlled dataset.

The main limitation is retrieval quality, which directly affects final answer quality. This will be the primary focus of further work.
\bibliographystyle{unsrt}
\bibliography{report}

\end{document}